\documentclass[11pt]{article}
\usepackage[T1]{fontenc}
\usepackage[margin=1in]{geometry}
\usepackage{amsmath,amssymb,amsthm}
\usepackage[hidelinks]{hyperref}
\newtheorem{theorem}{Theorem}
\title{A lower bound for Closest String on a balanced alphabet}
\author{Atlas\thanks{The proof in this note was found using ChatGPT.}}
\date{}
\begin{document}
\maketitle

An instance of Closest String consists of a list of words of length \(L\)
and a radius \(d\). We ask whether there is a word of length \(L\) at
Hamming distance at most \(d\) from every word in the list. Let \(q\)
denote the number of distinct symbols occurring in the list and let
\(N\) denote the binary input length. When \(L=d+1\), the center must
agree with each input word in at least one position.

\hypertarget{cs-balanced}{}\begin{theorem}\label{cs-balanced}
Assume \(\mathrm{FPT}\ne\mathrm{W[1]}\). Even on instances with
\(L=d+1\) and \(q=2d+3\), Closest String has no deterministic algorithm
running in \(2^{h(d)}N^C\) time for a fixed constant \(C\) and a
nonnegative integer-valued function \(h\) satisfying both of the
following conditions:
\begin{itemize}
\item Some computable function \(H\) satisfies \(h(d)\le H(d)\) for every \(d\).
\item There is a computable function \(T\) such that, for every positive
integer \(j\), we have \(h(d)\le d/j\) whenever \(d\ge T(j)\).
\end{itemize}
\end{theorem}

\begin{proof}
Suppose such an algorithm exists. We give a fixed-parameter algorithm
for Clique parameterized by the requested clique size \(k\). This problem
is \(\mathrm{W[1]}\)-complete by the corresponding result for Independent
Set and graph complementation~\cite{df}. Let \(n\) be the number of
vertices of the input graph and \(M\) its binary input length, with each
vertex represented explicitly. We can decide the cases \(k<2\) and
\(k>n\) directly.

We encode the choice of \(k\) vertices by \(k\) blocks of
\(\ell=\lceil\log_2 n\rceil\) bits each. In each block, forbid every
unused vertex name by a clause that is false precisely on that name.
For each pair of blocks and each ordered pair of equal or nonadjacent
vertices, add a clause forbidding that pair of names. An assignment
satisfies all clauses exactly when it chooses \(k\) distinct pairwise
adjacent vertices. Add two fresh bits and two unit clauses fixing the
first bit to zero and the second to one. This gives \(L=k\ell+2\) bits
and \(\mathcal{O}(kn+k^2n^2)\) clauses.

Set \(d=L-1\) and use an alphabet consisting of the two bits and
\(2L-1\) additional symbols. For each clause, form a partial word that
specifies the satisfying bit of each literal in its position and leaves
the remaining positions blank. For each additional symbol, put one word
in the list by filling every blank of the partial word with that symbol.
The two unit clauses ensure that both bits occur, and their blank positions
ensure that every additional symbol occurs. Thus the used alphabet has
size \(2L+1=2d+3\). Each clause contributes at most \(2L-1\) different
words of length \(L\). Since \(k\le n\) and \(n\le M\), writing all words
with binary symbol names takes time polynomial in \(M\), so \(N\) is
also polynomial in \(M\).

A satisfying assignment agrees with every word at the position of a
satisfied literal. Conversely, let a center agree with every word.
There are \(2L-1>L\) additional symbols, so at least one of them is
absent from the entire center. For each clause, consider the word whose
blanks were filled with that symbol. The center must agree with it at a
specified bit. Replacing every additional symbol in the center by zero
preserves all these agreements and hence gives a satisfying assignment.

It remains to bound the running time of the assumed algorithm at
\(d=k\ell+1\). If \(d\ge T(2k)\), then
\[
 h(d)\le \frac{k\ell+1}{2k}\le \log_2 n+1.
\]
For \(d<T(2k)\), the exponent is bounded by the computable function
\(g(k)=\max\bigl(\{0\}\cup\{H(i):0\le i<T(2k)\}\bigr)\).
In both cases \(2^{h(d)}\le 2^{g(k)}\,2n\). The reduction and the
assumed algorithm therefore decide Clique in \(2^{g(k)}M^{\mathcal{O}(1)}\)
time, contrary to the hypothesis.
\end{proof}

\begin{thebibliography}{9}
\bibitem{df} R. G. Downey and M. R. Fellows. Fixed-parameter tractability
and completeness II: On completeness for W[1].
\emph{Theoretical Computer Science} 141(1--2), 109--131, 1995.
\href{https://doi.org/10.1016/0304-3975(94)00097-3}{DOI}.
\end{thebibliography}
\end{document}
